\section{Claudeとの協働によるリーマン予想の証明主張}
Anthropicの数学者Levent Alpöge氏がClaudeとの協働でリーマン予想および一般化リーマン予想の証明をLean形式で公開したとする主張が広まった。Lean/Mathlibコミュニティが定理の一致とコンパイルを確認したとの報告もある。ただし形式的な検証の報告と数学的内容に対する独立した査読は区別する必要がある。Grokが示した主要X投稿は07:05頃で観測窓外に当たる。
\dailyxsource{2logicsによる共有（窓外）}{https://x.com/2logics/status/2124262966983827772}

\section{Inkling 2 Miniのローカル実行が話題に}
Thinking Machines Labの26.4Bパラメータのループ型MoEモデル「Inkling 2 Mini」が話題となった。約20GBでラップトップ動作可能とされ、Apache 2.0でオープンウェイト公開された。中間層の反復やNVFP4/MXFP8形式に加え、32GB Macでの実行報告、DeepSWEのスコアを巡る議論が見られた。リリース自体は11月20日で観測窓より前だが、技術者による検討が窓内で続いた。
\dailyxsource{Nikhilの関連投稿}{https://x.com/faltunik/status/2123989814651338877}

\section{米下院の中国製オープンモデル排除勧告}
米下院の委員会が、中国開発のオープンウェイトモデルを連邦システムや高リスクの重要インフラから排除するよう勧告したと報じられた。これを受け、Cursorを開発するAnysphereが対象モデルの段階的停止を検討しているとの投稿も共有された。下院報告書への言及はある一方、Cursorの公式声明はGrok収集時点で未確認であり、製品側の対応は確定情報として扱わない。
\dailyxsource{Jellycoolによる共有}{https://x.com/jellycoolgu/status/2124247617991516635}

\section{GPT-Live-1.1 APIとstandby料金}
OpenAIの全二重音声モデルGPT-Live-1.1のAPI提供が話題となった。X投稿では、アクティブ時が毎分0.05ドル、待機状態のstandbyモードが毎分0.02ドルと説明されている。常時接続の音声エージェントにおいて、会話していない時間の費用を下げる設計が注目された。公式発表自体は観測窓より前とされ、料金と仕様の厳密な確認は今後の課題である。
\dailyxsource{Holden.Chenの関連投稿}{https://x.com/jornyc/status/2124254746399961111}

\section{Metanoのエージェント実行時セキュリティ}
Metanoがステルス状態から公開に移り、1200万ドルのシード調達とAIエージェント向けランタイムセキュリティ・ガバナンス基盤の一般提供を発表したとの投稿が見られた。エージェントの行動をHiveグラフで可視化し、MCP対応や短時間導入をうたう。資金調達額や機能は投稿者の主張で、公式発表の直接確認は限定的である。主要X投稿は07:01頃で観測窓外となる。
\dailyxsource{SiMeltによる共有（窓外）}{https://x.com/zshjzls/status/2124261977400672686}

\section{Tencent Hy4のマルチモーダルモデル報告}
Tencentが総770Bパラメータ、アクティブ49BのMoEモデル「Hy4」をオープンウェイトで公開したとの情報が共有された。テキスト、画像、動画を扱い、100万トークンのコンテキスト長に対応するとの説明だ。ただし、これらはX投稿に列挙された仕様であり、Grokの収集では公式リポジトリやライセンスの確認には至っていない。利用可能性や実測性能は今後の確認事項である。
\dailyxsource{REJAULによる共有}{https://x.com/MRR1572/status/2124058246017130626}

\section{Transformers v5.24.0のモデル対応}
Hugging Face Transformers v5.24.0が、Inkling 2 Mini、Hy4-VL、NVIDIA Doryなどのネイティブサポートを追加したとの投稿があった。複数の新モデルを共通の実装基盤から利用できるようにする更新として注目される。ただし、対象モデルとリリース内容は投稿に基づくもので、GitHubのリリースノートは別途確認が必要である。主要投稿時刻は07:00頃で観測窓の境界に位置する。
\dailyxsource{2logicsの関連投稿}{https://x.com/2logics/status/2124261703382716827}

\section{METRがAI能力の時間地平を更新}
METRの2026年第4四半期の能力評価更新がX上で共有された。AIが研究開発関連のタスクで扱える仕事の長さ、すなわち時間地平が急速に延びているとの分析が中心で、評価そのものが難しくなっているとも指摘された。公式ブログの発表は観測窓前だが、窓内で測定方法や倍増ペースに関心が集まった。実務上の能力への外挿には、評価対象と測定条件の理解が必要になる。
\dailyxsource{YUK\_Kの関連投稿}{https://x.com/YUK_KND/status/2124249456295952402}

\section{医薬化学ベンチマークMedChemBench}
OpenAIの医薬化学向けベンチマーク「MedChemBench」が再び話題となった。99問から構成され、GPT-6.2 Astraが84.5\%を記録したとの投稿が共有された。専門領域でのモデル能力を測る試みとして注目されるが、数値は今回のX投稿に基づく報告として扱う。Grokが示した投稿は07:04頃で観測窓外に当たり、公式ページへの参照は背景情報となる。
\dailyxsource{Liam J. Wytcherleyによる共有（窓外）}{https://x.com/lwytch/status/2124262813376163937}

\section{AntigravityのGeminiクォータ改善}
GoogleのAntigravityについて、Geminiの利用クォータがリセットされ、キャッシュ改善によって22\%多くの作業が可能になったとの投稿が共有された。長時間のコーディングエージェント利用では、モデル性能だけでなくクォータやキャッシュ効率が実用性を左右する。開発者の投稿が引用されているものの、Grok収集時点でGoogle側の正式な仕様変更としての確認は限定的である。
\dailyxsource{Varun Mohanの関連投稿}{https://x.com/_mohansolo/status/2124070033534161407}

\section{Frostbiteコアの電力効率主張}
Architect LabsがAIを使って設計した「Frostbite」コアについて、従来CPUの4.59倍の性能／ワットをpre-silicon段階で達成したとの情報が共有された。AIによる半導体設計の可能性を示す主張として興味深いが、一次資料や評価条件はGrok収集時点で未確認である。実シリコン上の測定結果とも区別すべきであり、主要X投稿は07:03頃で観測窓外だった。
\dailyxsource{Artur Nowocinによる共有（窓外）}{https://x.com/arturnow/status/2124262408411586742}

\section{モデルが自身の重みを編集する研究}
オープンモデルが自身の重みを編集し、未見の問題で精度を向上させたとする研究が共有された。従来の推論時の文脈更新とは異なり、モデルのパラメータ自体を変更する方向性として注目される。投稿者はプロジェクトページを示しているが、独立した検証は未確認であり、性能向上の条件や再現性は今後の論文・コード公開を待つ必要がある。
\dailyxsource{Jiajun Fanの関連投稿}{https://x.com/FanJiajun67938/status/2124226655363481979}
